\chapter{Literature Review}

\section{Introduction}

The literature review examines how the research community has studied data corruption and model robustness, with emphasis on work that is directly relevant to individual and compound noise. The purpose is threefold: to establish what is already known, to identify the experimental settings that have been used, and to isolate the gap that the present project addresses.

The review was conducted using the IEEE Xplore Digital Library as the primary source, supplemented by ScienceDirect, SpringerLink, and Google Scholar for cross-checking metadata and for foundational papers. The search terms included combinations of \emph{noisy labels}, \emph{feature noise}, \emph{label noise robustness}, \emph{compound noise}, \emph{hybrid noise}, \emph{data corruption}, \emph{outliers}, \emph{missing values}, and \emph{classification}. Papers were included when they (i) appeared in a peer-reviewed journal or conference, (ii) addressed supervised classification under corrupted training data, and (iii) reported an experimental evaluation. Priority was given to four recent IEEE papers that collectively cover label noise, feature noise, hybrid noise, and compound noise. Foundational surveys and dataset papers were added to support the methodology.

The review is organised as follows. Section~2.2 describes the main families of existing approaches. Section~2.3 reviews the four core papers and the supporting literature in detail. Section~2.4 identifies the research gap, and Section~2.5 summarises the findings.

\section{Existing Systems}
Research on noisy data includes robust loss functions, sample selection and correction methods for label noise, as well as augmentation and robust estimation for feature corruption \cite{song2023survey,frenay2014classification,huber1964robust}. Classical classifiers remain useful for tabular data but are often evaluated under different datasets and noise assumptions \cite{hastie2009elements}.

\section{Review of Research Papers}
\paragraph{Fr\'enay and Verleysen.}
The survey distinguishes label errors that occur independently of an example from errors that depend on its true class or features \cite{frenay2014classification}. The project's label injector models uniform label noise by selecting a prescribed number of training rows at random and assigning each selected row a uniformly chosen alternative class.

\paragraph{Song et al.}
The survey groups deep-learning methods for noisy labels into robust architectures, regularisation, robust losses, loss adjustment, and sample selection \cite{song2023survey}. The present study uses classical classifiers with fixed settings and evaluates how their clean-test performance changes as training labels are corrupted.

\paragraph{Zhang et al.}
The proposed multilabel method uses relationships between labels and regularisation to address feature and label noise jointly \cite{zhang2020hybrid}. This motivates testing both input and label corruption in the same training split. The project evaluates single-label classifiers under individual and combined corruptions, with one target class per row.

\paragraph{Zeng et al.}
The authors show that appropriately chosen feature noise can improve generalisation for some deep networks trained with noisy labels \cite{zeng2025feature}. This motivates comparing label noise alone, Gaussian noise alone, and their combination in the official benchmark. Synergy measures departure from the sum of the individual drops. Whether Gaussian noise improved performance under label noise is assessed by comparing the combined condition directly with label noise alone.

\paragraph{Chen.}
The study addresses noisy labels and outliers together with a semisupervised generative model \cite{chen2023compound}. The present project includes outliers and label flips among four controlled corruption types and measures the resulting failure thresholds of classical tabular classifiers.

\section{Research Gap}

The reviewed literature establishes several important facts. Label noise harms generalisation, and many methods have been proposed to mitigate it \cite{song2023survey}. Feature noise can interact with label noise in non-obvious ways and can even improve robustness under some conditions \cite{zeng2025feature}. Hybrid and compound noise are recognised as realistic and difficult settings \cite{zhang2020hybrid,chen2023compound}. Classical models remain widely used on tabular data, but they are not inherently robust to corrupted labels \cite{frenay2014classification,song2023survey}.

At the same time, the literature leaves several gaps that are relevant to applied tabular machine learning.

\begin{enumerate}
  \item \textbf{Deep-learning bias.} The most comprehensive robustness studies focus on deep neural networks and image or graph benchmarks. Classical models on tabular data are underrepresented, even though they are common in medical, financial, and administrative applications.
  \item \textbf{Single-corruption bias.} Most studies examine label noise alone or feature noise alone. Compound noise is studied in selected settings, but there is no common protocol that enumerates all subsets of corruption types.
  \item \textbf{Limited severity sweep.} Robustness is usually reported at one or two noise rates. This makes it difficult to identify the point at which a model stops being useful.
  \item \textbf{No standard breaking-point measure.} The literature reports average accuracy or F1, but it does not define a model-specific failure threshold that can be compared across noise combinations.
  \item \textbf{Inconsistent evaluation protocols.} Datasets, splits, architectures, noise models, and metrics differ across papers, which limits direct comparison.
  \item \textbf{Weak treatment of synergy.} The combined effect of multiple noise types is often assumed to be additive or is discussed qualitatively, without an explicit decomposition into single-noise and interaction components.
\end{enumerate}

To the best of our knowledge, the reviewed literature does not report a single reproducible protocol that (i) enumerates all fifteen non-empty combinations of label noise, Gaussian feature noise, outliers, and missing values, (ii) sweeps five severity levels, (iii) evaluates four classical tabular models on both a binary and a multiclass dataset, and (iv) reports breaking points, aggregate robustness measures, and an explicit synergy analysis. This is the gap addressed by the present project.

The project does not claim to propose a new mitigation algorithm. Its contribution is a measurement framework and an empirical map. The framework makes the following aspects explicit and auditable:

\begin{itemize}
  \item the exact noise injectors and their parameters;
  \item the order in which compound noise is applied;
  \item the split and seed protocol;
  \item the definition of the breaking point;
  \item the aggregate robustness measures;
  \item the decomposition of compound effects into single-noise and interaction components;
  \item the provenance of every result row and the checksums of the final artefacts.
\end{itemize}

This positioning is important. The literature already provides strong evidence that noise harms models and that mitigation is possible. What it does not provide, in a single comparable protocol, is a map of \emph{where} classical tabular models break under individual and compound corruption. That map is the contribution of this project.

\section{Summary}

This chapter reviewed the main families of work on noisy data and model robustness. Label-noise research provides a mature taxonomy of robust architectures, regularisation methods, robust losses, loss adjustment, and sample selection \cite{song2023survey}. Feature-noise research shows that input corruption can interact with label corruption and can even improve generalisation in deep networks \cite{zeng2025feature}. Hybrid and compound-noise research shows that multiple corruption types must be considered together \cite{zhang2020hybrid,chen2023compound}. Classical robust learning provides the statistical and algorithmic background for handling outliers and noisy labels \cite{frenay2014classification,huber1964robust}.

The review also identified a clear gap: there is no systematic, reproducible, breaking-point-oriented comparison of classical tabular models under all combinations of label noise, Gaussian feature noise, outliers, and missing values. The next chapter describes the methodology designed to fill that gap.
